\subsubsection{Frontend}

\begin{figure}[H]
	\centering 
	\includegraphics[width=1.0\textwidth]{SpecificaTecnica/Autenticazione/UMLAutenticazioneBackend}
	\caption{Diagramma delle classi per Autenticazione e Registrazione: Backend}
\end{figure}

\subsubsubsection{CheckUserController}
\begin{figure}[H]
	\centering 
	\includegraphics[width=1\textwidth]{SpecificaTecnica/Autenticazione/CheckUserController}
	\caption{Autenticazione e Registrazione: CheckUserController}
\end{figure}

\paragraph{Descrizione}
\textit{CheckUserController} è il controller REST responsabile della gestione delle richieste HTTP di autenticazione.\\
Riceve le credenziali dell'utente, le delega al \textit{CheckUserService} e restituisce la risposta appropriata al client.

\paragraph{Attributi}
Questa classe non ha attributi propri.

\paragraph{Metodi e funzioni}
\begin{itemize}
	\item \textbf{create\_token(username: string): string} - attributo privato che rappresenta la funzione utilizzata per generare il token JWT a partire dall'username dell'utente autenticato.
	\item \textbf{login(request: Request, payload: UserSchema, service: \newline CheckUserServiceDep, response: Response): LoginResponse} - \newline metodo pubblico che gestisce la richiesta di login.\\
	Riceve i dati dell'utente tramite \textit{payload}, invoca il service per la verifica delle credenziali e restituisce una \textit{LoginResponse} contenente il token di accesso o un errore.
\end{itemize}